\section{哈密顿微分同胚群}
\stfhSubIndex{I152}{ha^mi^dun}{哈密顿（Hamiltonian）}{ha^mi^dun^wei^fen^tong^pei^qun}{哈密顿微分同胚群（Hamiltonian differomorphism group）}
\label{sec:hamiltonian-diffeomorphism-group}

由 $\omega$ 的闭性以及嘉当的“魔术”公式
\[
  \mathcal L_X\alpha=d(X\mathbin{\lrcorner}\alpha)
  +X\mathbin{\lrcorner}d\alpha,
\]
对一般向量场都有
$\mathcal L_X\omega=d(X\mathbin{\lrcorner}\omega)$。特别地，对 $M$ 上任意光滑函数 $f$，
\begin{equation}
  \mathcal L_{X_f}\omega=0.
  \label{eq:2.3.19}
\end{equation}

\begin{definition}
\label{def:symplectic-diffeomorphism-vector-field}
\stfhIndex{I267}{xin}{辛（symplectic）}
\stfhSubIndex{I268}{xin}{辛（symplectic）}{xin^wei^fen^tong^pei}{辛微分同胚（symplectic diffeomorphism）}
\stfhSubIndex{I269}{xin}{辛（symplectic）}{xin^xiang^liang^chang}{辛向量场（symplectic vector field）}
若微分同胚 $\phi:M\to M$ 满足 $\phi^*\omega=\omega$，则称 $\phi$ 是\emph{辛的}。以 $\Symp(M,\omega)$ 表示所有辛微分同胚的集合。若向量场 $X$ 满足 $\mathcal L_X\omega=0$，等价地，1-形式 $X\mathbin{\lrcorner}\omega$ 是闭的，则称 $X$ 是\emph{辛向量场}。
\end{definition}

显然，$\Symp(M,\omega)$ 构成 $\Diff(M)$ 的一个子群。事实上，由定义，任意辛微分同胚都保持由体积形式 $\omega^n/\omega!$ 诱导的刘维尔测度 $\mu_\omega$。

\begin{lemma}[刘维尔引理]
\label{lem:liouville-measure-preserved}
任意辛微分同胚 $\phi$ 都保持刘维尔测度；也就是说，对任意子集 $U\subset M$，
\[
  \mu_\omega(U)=\mu_\omega(\phi(U)).
\]
特别地，对任意辛向量场 $X$，
\[
  \mathcal L_X\omega^n=0.
\]
\end{lemma}

\subsection{\texorpdfstring{$\Ham(M,\omega)$ 的定义}{Ham(M, omega) 的定义}}

群 $\Symp(M,\omega)$ 的定义是自然的，因为它是辛形式 $\omega$ 的结构群。以 $\Symp_0(M,\omega)$ 表示 $\Symp(M,\omega)$ 的恒等连通分支。然而，$\Symp_0(M,\omega)$ 中还有另一个更为神秘的子群，称为\emph{哈密顿微分同胚群}。

先从哈密顿向量场的定义开始。

\begin{definition}
\label{def:hamiltonian-vector-field-isotopy}
\stfhSubIndex{I155}{ha^mi^dun}{哈密顿（Hamiltonian）}{ha^mi^dun^xiang^liang^chang}{哈密顿向量场（Hamiltonian vector field）}
\stfhSubIndex{I161}{ha^mi^dun^xiang^liang^chang}{哈密顿向量场（Hamiltonian vector field）}{ju^bu^ha^mi^dun^de}{局部哈密顿的（locally Hamiltonian）}
设 $(M,\omega)$ 是辛流形。
\begin{enumerate}[label=(\arabic*)]
  \item 若 $X\mathbin{\lrcorner}\omega$ 是正合的，则称向量场 $X$ 是\emph{哈密顿向量场}。
  \item 若微分同胚的同痕 $\phi^t$ 的生成向量场 $X_t$ 是\emph{辛的}（相应地，是\emph{哈密顿的}），则称 $\phi^t$ 是\emph{辛同痕}（相应地，\emph{哈密顿同痕}）。
\end{enumerate}
\end{definition}

以 $\operatorname{symp}(M,\omega)$（相应地，$\operatorname{ham}(M,\omega)$）表示辛向量场（相应地，哈密顿向量场）组成的子集。显然有包含关系
\[
  \operatorname{ham}(M,\omega)
  \subset\operatorname{symp}(M,\omega)
  \subset\mathfrak X(M).
\]
下面引理的证明留作习题。

\begin{lemma}
\label{lem:bracket-symplectic-fields-hamiltonian}
设 $X,Y$ 是两个辛向量场，则李括号 $[X,Y]$ 是与函数
\[
  f=-\omega(X,Y)
\]
相伴的哈密顿向量场。
\end{lemma}

\begin{proof}
证明留作习题。
\end{proof}

\begin{proposition}
\label{prop:symp-ham-lie-algebras}
\begin{enumerate}[label=(\arabic*)]
  \item $\operatorname{symp}(M,\omega)$ 与 $\operatorname{ham}(M,\omega)$ 都是 $\mathfrak X(M)$ 的李子代数。
  \item 有
  \[
    [\operatorname{symp}(M,\omega),\operatorname{symp}(M,\omega)]
    \subset\operatorname{ham}(M,\omega).
  \]
  特别地，
  \[
    [\operatorname{ham}(M,\omega),\operatorname{symp}(M,\omega)]
    \subset\operatorname{ham}(M,\omega),
  \]
  即 $\operatorname{ham}(M,\omega)$ 是 $\operatorname{symp}(M,\omega)$ 的正规李子代数。
  \item 有李代数的正合序列
  \[
    0\longrightarrow\operatorname{ham}(M,\omega)
    \longrightarrow\operatorname{symp}(M,\omega)
    \longrightarrow H^1(M,\mathbb R)\longrightarrow0,
  \]
  其中第一个映射是包含映射，第二个映射定义为
  \[
    X\longmapsto[X\mathbin{\lrcorner}\omega],
  \]
  而 $H^1(M,\mathbb R)$ 上取平凡括号。
\end{enumerate}
\end{proposition}

\begin{proof}
(1) 与 (2) 是引理 \ref{lem:bracket-symplectic-fields-hamiltonian} 的直接推论。(3) 也很容易由定义验证，细节留给读者。
\end{proof}

$\Symp(M,\omega)$ 是几何结构——辛结构 $\omega$——的自同构群，而 $\operatorname{symp}(M,\omega)$ 是与之相伴的李代数。当 $M$ 为闭流形时，我们在 $\Symp(M,\omega)$ 上赋予由 $\Diff(M)$ 的 $C^\infty$ 拓扑诱导的拓扑。

命题 \ref{prop:symp-ham-lie-algebras}(2) 提示，李群 $\Symp(M,\omega)$ 中应该存在一个正规子群，其相应的李代数是 $\operatorname{ham}(M,\omega)$。事实上，这样的子群确实存在；但它的定义不如 $\Symp(M,\omega)$ 的定义自然，特别是当我们试图在它上面赋予拓扑时更是如此。

设
$H:[0,1]\times M\to\mathbb R$ 是随时间变化的光滑函数族。记
$H_t=H(t,\mathord\cdot):M\to\mathbb R$，并以 $\phi_H^t$ 表示常微分方程
\[
  \dot x=X_{H_t}(x)
\]
的流。为简化记号，也写 $X_{H_t}(x)=X_H(t,x)$，把 $X_H$ 看作时变向量场。

\begin{definition}
\label{def:hamiltonian-diffeomorphism}
若存在时变哈密顿函数
\[
  H:[0,1]\times M\longrightarrow\mathbb R
\]
使微分同胚 $\phi:M\to M$ 满足 $\phi=\phi_H^1$ 且 $\phi_H^0=\operatorname{id}$，则称 $\phi$ 是\emph{哈密顿微分同胚}。在这种情形下记作 $H\mapsto\phi$。所有哈密顿微分同胚组成的集合记为 $\Ham(M,\omega)$。
\stfhSubIndex{I151}{ha^mi^dun}{哈密顿（Hamiltonian）}{ha^mi^dun^wei^fen^tong^pei}{哈密顿微分同胚（Hamiltonian diffeomorphism）}
\end{definition}

由定义并不能直接看出 $\Ham(M,\omega)$ 构成群，虽然它显然是 $\Symp(M,\omega)$ 的子集。这个事实的证明涉及一种有趣的哈密顿函数代数；它由霍费尔 (H90) 引入，并将在研究哈密顿微分同胚的几何以及弗洛尔同调理论时发挥重要作用。先从下面的引理开始。

\begin{lemma}
\label{lem:symplectic-pullback-hamiltonian-field}
设 $h:M\to\mathbb R$ 是函数，$X_h$ 是与之相伴的哈密顿向量场。则对任意辛微分同胚 $\psi:M\to M$，有
\[
  \psi^*(X_h)=X_{h\circ\psi}.
\]
\end{lemma}

\begin{proof}
我们将证明
\begin{equation}
  \psi^*(X_h)\mathbin{\lrcorner}\omega=d(h\circ\psi).
  \label{eq:2.3.20}
\end{equation}
由于 $\psi$ 是辛的，$\psi^*\omega=\omega$，所以
\[
  \psi^*(X_h)\mathbin{\lrcorner}\omega
  =\psi^*(X_h)\mathbin{\lrcorner}\psi^*\omega.
\]
另一方面，右边正是
\[
  \psi^*(X_h\mathbin{\lrcorner}\omega)
  =\psi^*(dh)=d(h\circ\psi).
\]
合并两式即得 \eqref{eq:2.3.20}。
\end{proof}

\begin{proposition}
\label{prop:hamiltonian-composition-inverse}
设 $\phi,\psi$ 是两个哈密顿微分同胚，$H,K:[0,1]\times M\to\mathbb R$ 是它们的生成哈密顿量。则：
\begin{enumerate}[label=(\arabic*)]
  \item $\overline H\mapsto\phi^{-1}$，其中 $\overline H$ 定义为
  \begin{equation}
    \overline H(t,x):=-H(t,\phi_H^t(x)).
    \label{eq:2.3.21}
  \end{equation}
  \item $H\#K\mapsto\phi\psi$，其中 $H\#K:[0,1]\times M\to\mathbb R$ 定义为
  \begin{equation}
    H\#K(t,x):=H(t,x)+K(t,(\phi_H^t)^{-1}(x)).
    \label{eq:2.3.22}
  \end{equation}
\end{enumerate}
特别地，$\Ham(M,\omega)$ 是 $\Symp(M,\omega)$ 的子群。
\end{proposition}

\begin{proof}
设 $Y_t$ 是生成同痕 $(\phi_H^t)^{-1}$ 的向量场，即
\[
  Y_t=\frac d{dt}\bigl((\phi_H^t)^{-1}\bigr)\circ\phi_H^t.
\]
计算得
\begin{align*}
 0&=\frac d{dt}(\operatorname{id})(x)
   =\frac d{dt}\bigl((\phi_H^t)^{-1}\bigr)\circ\phi_H^t(x)\\
  &=d(\phi_H^t)^{-1}\left(\frac{d\phi_H^t}{dt}(x)\right)
    +Y_t\bigl((\phi_H^t)^{-1}\circ\phi_H^t(x)\bigr).
\end{align*}
因此
\begin{align*}
 Y_t(x)
 &=-d(\phi_H^t)^{-1}\left(\frac{d\phi_H^t}{dt}(x)\right)\\
 &=-d(\phi_H^t)^{-1}\bigl(X_{H_t}(\phi_H^t(x))\bigr)\\
 &=-(\phi_H^t)^*X_{H_t}(x)
   =X_{-H_t\circ\phi_H^t}(x).
\end{align*}
这就证明了 (1)。

类似地，对复合同痕 $t\mapsto\phi_H^t\circ\phi_K^t$ 求导，作一个直接计算即可证明 (2)。细节留给读者作为习题。
\end{proof}

\begin{exercise}
\label{ex:composition-hamiltonian}
证明函数 \eqref{eq:2.3.22} 生成同痕 $\phi_H^t\circ\phi_K^t$。
\end{exercise}

现在要证明，$\Ham(M,\omega)$ 恰是 $\Symp(M,\omega)$ 中以 $\operatorname{ham}(M,\omega)$ 为相伴李代数的李子群。换言之，需要证明 $\Ham(M,\omega)$ 的任意单参数子群 $\{\phi^s\}$ 都形如 $\phi_h^s$，即它是某个函数 $h\in C^\infty(M)$ 所对应的\emph{自治}哈密顿向量场之流。事实表明，在着手证明之前，需要先讨论 $\Ham(M,\omega)$ 的拓扑；证明将用到最初由奥古斯丁·巴尼亚加（Augustin Banyaga）(Ba78) 给出的一些非平凡论证。
\stfhIndex{I059}{ba^ni^ya^jia}{巴尼亚加（Banyaga）}

\subsection{\texorpdfstring{$\Ham(M,\omega)$ 的一般拓扑}{Ham(M, omega) 的一般拓扑}}

先简要讨论由 $\Diff(M)$ 的 $C^\infty$ 拓扑诱导出的 $\Symp(M,\omega)$ 上的 $C^\infty$ 拓扑。下述局部可缩性结论本质上是温斯坦 (Wn73) 关于拉格朗日子流形的达布邻域定理的推论；我们将在第 3.3 节给出证明，并会在后续章节详细讨论拉格朗日子流形。

\begin{proposition}
\label{prop:symp-topology}
关于 $C^\infty$ 拓扑，群 $\Symp(M,\omega)$ 是 $\Diff(M)$ 的一个闭且局部闭子群。由 $\Diff(M)$ 诱导出的 $\Symp(M,\omega)$ 上的拓扑是局部可缩的。
\end{proposition}

类似地，通过包含映射 $\Ham(M,\omega)\subset\Diff(M)$ 所诱导的拓扑来定义 $\Ham(M,\omega)$ 的 $C^\infty$ 拓扑。然而，与 $\Symp(M,\omega)$ 的情形不同，$\Ham(M,\omega)$ 是否为 $\Diff(M)$ 或 $\Symp(M,\omega)$ 的闭子群，是一个高度非平凡的问题。这个问题与所谓的 $C^\infty$ \emph{通量猜想} (Ba78) 密切相关；小野薫 (On06) 已在完全一般的情形下证明了该猜想。下一节将进一步讨论通量同态。

现在回到给定的单参数子群 $\phi^s\in\Ham(M,\omega)$。

\begin{remark}
\label{rem:one-parameter-subgroup-smoothness}
回忆，对\emph{有限维}李群，任意单参数子群都会自动定义该李群上微分同胚的光滑同痕。然而，由于 $\Ham(M,\omega)$ 是无限维群，一个单参数子群是否会在 $\Ham(M,\omega)$ 上定义光滑（甚至 $C^1$）路径并不清楚；也就是说，并不清楚曲线 $s\mapsto\phi^s(x)$ 是否是 $M$ 上的可微路径。
\end{remark}

鉴于此，我们总是考虑 $\Ham(M,\omega)$ 中的光滑族 $\{\phi^s\}$，这里“光滑”是指它作为 $\Diff(M)$ 中的族是光滑的。于是一个自然问题是：这样的同痕能否由哈密顿流得到？也就是说，是否存在时变哈密顿量
$K:[0,1]\times M\to\mathbb R$，使
\begin{equation}
  \phi^s=\phi_K^s\circ\phi^0.
  \label{eq:2.3.23}
\end{equation}
给这样的族一个明确名称会很有用。

\begin{definition}
\label{def:hamiltonian-path}
一个\emph{（光滑）哈密顿路径} $\lambda:[0,1]\to\Ham(M,\omega)$，是由光滑映射
\[
  \Lambda:[0,1]\times M\longrightarrow M
\]
通过 $\lambda(t)(x):=\Lambda(t,x)$ 定义的映射，并满足：
\begin{enumerate}[label=(\arabic*)]
  \item 其导数
  \[
    \dot\lambda(t)=\frac{\partial\lambda}{\partial t}
      \circ(\lambda(t))^{-1}
  \]
  是哈密顿的；也就是说，对所有 $t\in[0,1]$，1-形式
  $\dot\lambda(t)\mathbin{\lrcorner}\omega$ 都是正合的。若函数
  $H:\mathbb R\times M\to\mathbb R$ 满足
  \[
    \lambda(t)=\phi_H^t(\lambda(0)),
    \qquad\text{等价地}\qquad
    dH_t=\dot\lambda(t)\mathbin{\lrcorner}\omega,
  \]
  则称 $H$ 为 $\lambda$ 的\emph{生成哈密顿量}。
  \stfhSubIndex{I150}{ha^mi^dun}{哈密顿（Hamiltonian）}{sheng^cheng^ha^mi^dun^liang}{生成哈密顿量（generating Hamiltonian）}
  \item 微分同胚
  $\lambda(0):=\Lambda(0,\mathord\cdot):M\to M$ 是哈密顿微分同胚。
\end{enumerate}
\end{definition}

以 $\mathcal P^{\mathrm{ham}}(\Symp(M,\omega))$ 表示所有哈密顿路径
$\lambda:[0,1]\to\Symp(M,\omega)$ 的集合，以
$\mathcal P^{\mathrm{ham}}(\Symp(M,\omega),\operatorname{id})$ 表示其中满足 $\lambda(0)=\operatorname{id}$ 的路径集合。用上述相应映射 $\Lambda$ 的 $C^\infty$ 拓扑，在这两个空间上赋予显然的诱导拓扑。

引入映射 $H\mapsto\phi_H$ 的逆映射常常很有用。先回忆，常值函数生成常值路径；因此，为了唯一定义其逆，必须作适当的归一化。

\begin{definition}
\label{def:mean-normalized}
设 $(M,\omega)$ 是连通闭辛流形。若函数 $h:M\to\mathbb R$ 满足
\begin{equation}
  \int_M h\omega^n=0,
  \label{eq:2.3.24}
\end{equation}
则称 $h$ 是\emph{均值归一化的}。以 $C_m^\infty(M)$ 表示所有均值归一化光滑函数的集合。
\end{definition}

\begin{remark}
\label{rem:normalization-hamiltonian}
在把哈密顿向量场同其相伴的哈密顿函数联系起来时，我们经常需要消除常值函数带来的歧义。尤其是，当试图把由时变哈密顿函数构造的不变量解释为相应哈密顿路径的不变量时，更是如此。这种归一化会进入辛不变量的构造：既会出现在作用量泛函的临界点理论研究中，也会出现在哈密顿纤维化研究中。讨论哈密顿周期轨道的弗洛尔同调及其相伴的谱不变量时，我们还会回到这个归一化问题。
\end{remark}

下面的定义引自 (OhM07)。

\begin{definition}
\label{def:tangent-developing-maps}
设 $\lambda\in\mathcal P^{\mathrm{ham}}(\Symp(M,\omega),\operatorname{id})$，$H$ 是生成给定哈密顿路径 $\lambda$ 的归一化哈密顿量，即 $\lambda=\phi_H$。定义两个映射
\[
 \operatorname{Tan},\operatorname{Dev}:
 \mathcal P^{\mathrm{ham}}(\Symp(M,\omega),\operatorname{id})
 \longrightarrow C^\infty([0,1]\times M,\mathbb R)
\]
如下：
\[
 \operatorname{Tan}(\lambda)(t,x):=H(t,\phi_H^t(x)),
 \qquad
 \operatorname{Dev}(\lambda)(t,x):=H(t,x).
\]
称它们分别为\emph{切向映射}和\emph{展开映射}。切向映射的像称为 $\lambda$（或 $H$）的\emph{卷绕哈密顿量}。
\stfhIndex{I239}{juan^rao^ha^mi^dun^liang}{卷绕哈密顿量（rolled Hamiltonian）}
\stfhIndex{I284}{qie^xiang^ying^she}{切向映射（tangent map）}
\stfhIndex{I100}{zhan^kai^ying^she}{展开映射（developing map）}
\end{definition}

切向映射对应于路径切向量的映射。把通常的生成哈密顿量 $H$ 赋给一条哈密顿路径，对应于李群理论中的展开映射：先用切向映射，再作右平移，就可以把李群上的任意可微路径“展开”为其李代数中的路径。

我们将在第 2.4 节证明巴尼亚加 (Ba78) 的如下基本定理。

\begin{theorem}[巴尼亚加]
\label{thm:banyaga-smooth-ham-path}
$\Diff(M)$ 中由哈密顿微分同胚组成的任意光滑路径都是哈密顿路径。
\end{theorem}

暂且假设这个定理，回到 $\Ham(M,\omega)\subset\Diff(M)$ 中光滑单参数子群 $\phi^s$ 的研究。由该定理，$\phi^s$ 是哈密顿路径；也就是说，$X\mathbin{\lrcorner}\omega$ 是正合的，因而对某个光滑函数 $h:M\to\mathbb R$ 有 $X\mathbin{\lrcorner}\omega=dh$。上述讨论可总结如下。

\begin{corollary}
\label{cor:lie-algebra-ham}
$\Ham(M,\omega)$ 的李代数是由哈密顿向量场组成的集合 $\operatorname{ham}(M,\omega)$。
\end{corollary}

\begin{proof}
设 $\phi^s$ 是 $\Ham(M,\omega)\subset\Diff(M)$ 中的光滑单参数子群。由群性质，生成向量场
\[
  X_s:=\frac{d\phi^s}{ds}\circ(\phi^s)^{-1}
\]
与 $s$ 无关。把这个公共向量场记为 $X=X_0$。现在需要证明 $X$ 确为哈密顿的，即 $X\mathbin{\lrcorner}\omega$ 是正合的。但由定理 \ref{thm:banyaga-smooth-ham-path}，$\phi^s$ 是哈密顿路径，所以 $X\mathbin{\lrcorner}\omega$ 是正合的；即存在光滑函数 $h:M\to\mathbb R$，使 $X\mathbin{\lrcorner}\omega=dh$。证明完毕。
\end{proof}

回忆引理 \ref{lem:bracket-symplectic-fields-hamiltonian} 中的恒等式
$[X_f,X_g]=-X_{\{f,g\}}$。它说明赋值
\[
  f\in C^\infty(M)/\mathbb R\longmapsto-X_f
\]
定义李代数同构，故有如下结论。

\begin{proposition}
\label{prop:functions-lie-algebra-ham}
关于由 $\omega$ 相伴的泊松括号，向量空间 $C^\infty(M)/\mathbb R$ 与 $\Ham(M,\omega)$ 的李代数 $\operatorname{ham}(M,\omega)$ 典范同构。
\end{proposition}

现在从一般李群理论的观点重新表述上述讨论。首先有自然认同：
\begin{enumerate}[label=(\arabic*)]
  \item $T_{\operatorname{id}}\Ham(M,\omega)\cong$哈密顿向量场的集合；
  \item 哈密顿向量场的集合 $\cong$ 正合 1-形式的集合；
  \item $C_m^\infty(M)\cong$ 正合 1-形式的集合。
\end{enumerate}
其次，由\emph{李群理论的一般事实}，左乘（复合）
\begin{equation}
  L_\psi:\Ham(M,\omega)\longrightarrow\Ham(M,\omega);
  \qquad\phi\longmapsto\psi\circ\phi
  \label{eq:2.3.25}
\end{equation}
的微分
\[
 dL_\psi:T_{\operatorname{id}}\Ham(M,\omega)
 \longrightarrow T_\psi\Ham(M,\omega)
\]
给出 $T_{\operatorname{id}}\Ham(M,\omega)$ 与 $T_\psi\Ham(M,\omega)$ 之间的同构。注意，辛微分同胚 $\psi$ 在 $\Ham(M,\omega)$ 上的共轭
$\phi\mapsto\psi\phi\psi^{-1}$，是李群 $\Ham(M,\omega)$ 在其李代数——\emph{哈密顿向量场集合}——上的\emph{伴随作用}的一个特例。用 $C_m^\infty(M)$ 中的函数表示时，诱导的伴随作用为
$h\mapsto h\circ\psi^{-1}$。

对任意辛流形，$C^\infty(M)/\mathbb R$ 上有一个自然的伴随不变范数，定义为
\[
  \operatorname{osc}(h):=\max_{x\in M}h(x)-\min_{x\in M}h.
\]
李群的一般事实告诉我们，$\operatorname{osc}$ 在 $T\Ham(M,\omega)$ 上诱导一个双不变芬斯勒度量。这恰是 $\Ham(M,\omega)$ 上的霍费尔芬斯勒度量。对于这个度量，有恒等式
\[
 \int_0^1\|\dot\phi_t\|_{\phi_t}\,dt
 :=\int_0^1\|(d\phi_t)^{-1}\cdot\dot\phi_t\|_{\operatorname{id}}\,dt
 =\int_0^1\operatorname{osc}(H_t)\,dt,
\]
其中 $\phi_t=\phi_H^t$，而 $H:M\times[0,1]\to\mathbb R$ 是生成同痕 $\phi_t$ 的归一化哈密顿量。霍费尔 (H90) 的惊人洞见是：这个不变度量的几何与动力学，会为辛拓扑和哈密顿动力学中直至 $C^0$ 拓扑层次的问题提供丰富来源。
